\documentclass[12pt]{article}
\usepackage[utf8]{inputenc}
\usepackage[T1]{fontenc}
\usepackage[margin=1in]{geometry}
\usepackage{setspace}
\usepackage{xcolor}
\usepackage[colorlinks=true, urlcolor=blue!60!black, linkcolor=blue!60!black]{hyperref}
\setlength{\parindent}{0em}
\setlength{\parskip}{0.6em}
\renewcommand{\baselinestretch}{1.0}
\begin{document}
\begin{center}
{\LARGE \textit{United States v. Pink}, 315 U.S. 203 (1942)\par}
\vspace{0.5em}
{\large Case Brief --- Two-Page Edition\par}
\end{center}

\section{Memory Jogger}
A sole executive agreement made incident to recognizing the Soviet Union overrode New York law and gave the United States a nationalized Russian insurer's New York assets (\href{https://supreme.justia.com/cases/federal/us/315/203/}{315 U.S.\ 203}).

\section{Detailed Case Facts}
New York liquidated the New York branch of the First Russian Insurance Company, which Soviet decrees had nationalized, and planned to pay a surplus of over \$1 million to foreign creditors. When President Roosevelt recognized the Soviet Union in 1933, the Litvinov Assignment, made without the Senate, assigned Soviet claims against Americans to the United States.

\section{Procedural History}
The United States sued Superintendent Pink for the surplus. Relying on \textit{Moscow Fire}, the New York courts dismissed, applying New York's policy against confiscation.

\section{Judicial Votes}
Reversed, 5--2. Douglas wrote for Black, Frankfurter, Murphy, and Byrnes; Frankfurter concurred; Stone, joined by Roberts, dissented; Reed and Jackson did not participate.

\section{Holding}
Extending \textit{Belmont}, the President's recognition power includes the power to settle claims by sole executive agreement, which has the same supremacy over state law and policy as a treaty. Title to the surplus passed to the United States, and the Fifth Amendment did not bar preferring it over foreign creditors.

\section{Analysis of Opinions}
Douglas held that recognition and claims settlement were one bargain and that removing obstacles to recognition was a ``modest implied power'' of the President as ``sole organ'' in foreign affairs; letting New York refuse enforcement would revive the friction the President had resolved. Frankfurter warned against applying domestic concepts like ``situs'' and ``comity'' to international disputes. Stone argued that recognition does not reach New York assets and that the Assignment transferred only rights the Soviets actually had.

\section{Examples: Future Applications}
\textbf{Same-Side:} a recognition agreement assigning a foreign airline's U.S.\ assets preempts state rules. \textbf{Opposite-Side:} an agreement overriding a federal sanctions statute likely fails under \textit{Youngstown} and \textit{Medell\'in}. \textbf{Fence-Sitter:} a settlement extinguishing American creditors' attachments raises takings questions.

\section{Critique}
The Supremacy Clause does not list executive agreements, and Stone warned that the Executive could bypass the Senate to rewrite domestic private law; \textit{Medell\'in} later narrowed \textit{Pink}.

\section{Key Quotations}
\begin{quote}
``Power to remove such obstacles to full recognition as settlement of claims of our nationals ... certainly is a modest implied power of the President who is the `sole organ of the federal government in the field of international relations.'\,''
\end{quote}
\noindent --- Justice Douglas, for the Court.

\end{document}
